% MAC0216 - Técnicas de Programação I
% Trabalho Prático 2 - Biblioteca Modular em C
% Relatório Técnico


\documentclass[12pt,a4paper]{article}


\usepackage[T1]{fontenc}
\usepackage[utf8]{inputenc}
\usepackage[brazilian]{babel}
\usepackage{indentfirst}
\usepackage{csquotes}


\usepackage[a4paper,margin=2.5cm]{geometry}

\usepackage{setspace}
\onehalfspacing

\usepackage{parskip}
\setlength{\parindent}{1.25cm}

\usepackage{microtype}


\usepackage{lmodern}
\usepackage{graphicx}
\usepackage{xcolor}
\usepackage{booktabs}
\usepackage{array}
\usepackage{enumitem}
\usepackage{float}
\usepackage[font=small,labelfont=bf]{caption}
\usepackage{xspace}

\graphicspath{{figuras/}}


\usepackage[hidelinks]{hyperref}
\usepackage{url}

\hypersetup{
  pdftitle={Biblioteca Modular em C - Relatório Técnico},
  pdfsubject={MAC0216 - Técnicas de Programação I - Trabalho Prático 2},
  pdfkeywords={C, biblioteca, Makefile, ligação estática, ligação dinâmica,
    testes, Git, aprendizado de máquina}
}


\usepackage{listings}

\definecolor{codeback}{RGB}{248,248,248}
\definecolor{codeframe}{RGB}{215,215,215}
\definecolor{codekey}{RGB}{0,80,160}
\definecolor{codecomment}{RGB}{100,120,100}
\definecolor{codestring}{RGB}{160,50,30}
\definecolor{codeline}{RGB}{150,150,150}


\lstset{
  extendedchars=true,
  literate=%
    {á}{{\'a}}1 {à}{{\`a}}1 {â}{{\^a}}1 {ã}{{\~a}}1 {ä}{{\"a}}1
    {é}{{\'e}}1 {ê}{{\^e}}1 {í}{{\'i}}1 {ó}{{\'o}}1 {ô}{{\^o}}1
    {õ}{{\~o}}1 {ú}{{\'u}}1 {ü}{{\"u}}1 {ç}{{\c c}}1
    {Á}{{\'A}}1 {À}{{\`A}}1 {Â}{{\^A}}1 {Ã}{{\~A}}1
    {É}{{\'E}}1 {Ê}{{\^E}}1 {Í}{{\'I}}1 {Ó}{{\'O}}1 {Ô}{{\^O}}1
    {Õ}{{\~O}}1 {Ú}{{\'U}}1 {Ç}{{\c C}}1
    {°}{{\textdegree}}1
}

% Estilo para código-fonte em C (com numeração de linhas).
\lstdefinestyle{codigo}{
  language=C,
  backgroundcolor=\color{codeback},
  basicstyle=\ttfamily\footnotesize,
  keywordstyle=\color{codekey}\bfseries,
  commentstyle=\color{codecomment}\itshape,
  stringstyle=\color{codestring},
  numbers=left,
  numberstyle=\tiny\color{codeline},
  numbersep=8pt,
  frame=single,
  rulecolor=\color{codeframe},
  framesep=5pt,
  breaklines=true,
  breakatwhitespace=false,
  postbreak=\mbox{\textcolor{codeline}{$\hookrightarrow$}\space},
  showstringspaces=false,
  tabsize=4,
  captionpos=b,
  xleftmargin=15pt,
  aboveskip=1em,
  belowskip=1em
}

% Estilo para trechos do Makefile.
\lstdefinestyle{makefile}{
  style=codigo,
  language={[gnu]make}
}

% Estilo para sessões de terminal (sem numeração).
\lstdefinestyle{terminal}{
  language={},
  backgroundcolor=\color{codeback},
  basicstyle=\ttfamily\footnotesize,
  numbers=none,
  frame=single,
  rulecolor=\color{codeframe},
  framesep=5pt,
  breaklines=true,
  showstringspaces=false,
  tabsize=4,
  captionpos=b,
  xleftmargin=0pt,
  aboveskip=1em,
  belowskip=1em
}

\lstset{style=codigo}

\renewcommand{\lstlistingname}{Listagem}
\renewcommand{\lstlistlistingname}{Lista de Listagens}


% Marca o que ainda falta escrever (remover todas as ocorrências antes da entrega).
\newcommand{\preencher}[1]{\textcolor{red}{[#1]}}

% TODO: trocar pelo nome escolhido para a biblioteca.
\newcommand{\biblioteca}{\texttt{\preencher{Biblioteca de ML}}\xspace}
\newcommand{\cmd}[1]{\texttt{#1}}
\newcommand{\opt}[1]{\texttt{-{}-#1}}


\begin{document}

\hypersetup{pageanchor=false}

%------------------------------------------------------------------------------
% Capa
%------------------------------------------------------------------------------
\begin{titlepage}
\centering

{\large UNIVERSIDADE DE SÃO PAULO\par}
{\large Instituto de Matemática, Estatística e Ciência da Computação\par}
\vspace{0.3cm}
{\large MAC0216 --- Técnicas de Programação I\par}

\vfill

{\Huge\bfseries \biblioteca\par}
\vspace{0.6cm}
{\Large Biblioteca Modular em C para Aprendizado de Máquina\par}
\vspace{0.4cm}
{\large Relatório Técnico --- Trabalho Prático 2\par}

\vfill

{\large\bfseries Integrantes\par}
\vspace{0.4cm}
\begin{tabular}{l@{\hspace{1.5cm}}l}
  Felipe Brancalhão Curcio  & 16985106 \\ [4pt]
  Jaime Cavalcanti Guedes   & 16875204 \\ [4pt]
  João Pedro Silva Amadio   & 13672533 \\ [4pt]
  Murilo Santos de Castro   & 17015895 \\ [4pt]
\end{tabular}

\vfill

{\large São Paulo\par}
{\large \the\year\par}

\end{titlepage}


\pagenumbering{roman}
\tableofcontents
\newpage
\pagenumbering{arabic}
\hypersetup{pageanchor=true}

%==============================================================================
\section{Introdução}\label{sec:introducao}
%==============================================================================
% Enunciado, Seção 7, item 1: descrição do problema e do domínio escolhido.

Este relatório descreve o projeto e a implementação da \biblioteca, uma
biblioteca modular escrita em linguagem C e voltada para a área de aprendizado
de máquina (\emph{Machine Learning}), acompanhada de uma aplicação cliente de
linha de comando. O código-fonte está disponível no repositório
\href{https://github.com/Jaime-Cg-si/MAC0216-TP2-Bibliotecas-Make.git}{TP2}.

\end{document}
